\documentclass{article}
%packages
\usepackage{fullpage}
\usepackage{xspace}
\usepackage[dvipsnames]{xcolor}
\usepackage{amsmath, amsthm, amsfonts, amssymb}
\usepackage{stmaryrd}
\usepackage{float}
\usepackage{mathtools}
\usepackage{mathrsfs}
\usepackage[shortlabels]{enumitem}
\usepackage{stackrel}
\usepackage[framemethod = TikZ]{mdframed}

%theorems
\newtheorem{theorem}{Teorema}[subsection]

\newtheorem{algorithm}{Algorithm}
\newtheorem{axiom}{Axioma}
\newtheorem{case}{Case}
\newtheorem{claim}{Afirmación}
\newtheorem{conclusion}{Conclusión}
\newtheorem{condition}{Condition}
\newtheorem{conjecture}[theorem]{Conjetura}
\newtheorem{corollary}[theorem]{Corolario}
\newtheorem{criterion}{Criterion}
\newtheorem{proposition}[theorem]{Proposición}
\newtheorem{lemma}[theorem]{Lema}

\theoremstyle{definition}
\newtheorem{definition}[theorem]{Definición}

\newtheorem{exercise}{Ejercicio}

\newtheorem*{notation}{Notación}
\newtheorem{problem}{Problema}[section]
\newtheorem*{obs}{Observación}
\newtheorem*{addendum}{Addendum}

\newtheorem*{tarea}{\wea{Tarea}}

\theoremstyle{remark}
\newtheorem{remark}[theorem]{Remark}
\newtheorem*{solution}{\textbf{Solución}}
\newtheorem{summary}{Summary}
\newtheorem{example}[theorem]{\textbf{Ejemplo}}

\numberwithin{equation}{section}

\renewcommand*{\proofname}{\textbf{Demostración}}

\surroundwithmdframed[
  topline=false,
  rightline=false,
  bottomline=false,
  leftmargin=\parindent,
  skipabove=\medskipamount,
  skipbelow=\medskipamount
]{proof}

\surroundwithmdframed[
  topline=false,
  rightline=false,
  bottomline=false,
  leftmargin=\parindent,
  skipabove=\medskipamount,
  skipbelow=\medskipamount
]{solution}

%commands
\newcommand{\A}{\mathcal{A}}
\newcommand{\B}{\mathcal{B}}
\newcommand{\C}{\mathbb{C}}
\newcommand{\D}{\mathbb{D}}
\newcommand{\F}{\mathbb{F}}
\newcommand{\N}{\mathbb{N}}
\renewcommand{\O}{\mathcal{O}}
\renewcommand{\P}{\mathcal{P}}
\newcommand{\Q}{\mathbb{Q}}
\newcommand{\R}{\mathbb{R}}
\renewcommand{\S}{\mathbb{S}}
\newcommand{\U}{\mathcal{U}}
\newcommand{\Z}{\mathbb{Z}}

\newcommand{\Acal}{\mathcal{A}}
\newcommand{\Bcal}{\mathcal{B}}
\newcommand{\Ccal}{\mathcal{C}}
\newcommand{\Dcal}{\mathcal{D}}
\newcommand{\Ecal}{\mathcal{E}}
\newcommand{\Fcal}{\mathcal{F}}
\newcommand{\Gcal}{\mathcal{G}}
\newcommand{\Hcal}{\mathcal{H}}
\newcommand{\Ical}{\mathcal{I}}
\newcommand{\Jcal}{\mathcal{J}}
\newcommand{\Kcal}{\mathcal{K}}
\newcommand{\Lcal}{\mathcal{L}}
\newcommand{\Mcal}{\mathcal{M}}
\newcommand{\Ncal}{\mathcal{N}}
\newcommand{\Qcal}{\mathcal{Q}}
\newcommand{\Rcal}{\mathcal{R}}
\newcommand{\Scal}{\mathcal{S}}
\newcommand{\Tcal}{\mathcal{T}}
\newcommand{\Ucal}{\mathcal{U}}
\newcommand{\Vcal}{\mathcal{V}}
\newcommand{\Wcal}{\mathcal{W}}
\newcommand{\Xcal}{\mathcal{X}}
\newcommand{\Ycal}{\mathcal{Y}}
\newcommand{\Zcal}{\mathcal{Z}}

\renewcommand{\a}{\alpha}
\renewcommand{\b}{\beta}
\renewcommand{\d}{\mathrm{d}}
\newcommand{\e}{\varepsilon}
\renewcommand{\t}{\tau}
\newcommand{\lam}{\lambda}
\newcommand{\Lam}{\Lambda}

\newcommand{\vphi}{\varphi}
\newcommand{\oo}{\infty}
\newcommand{\longto}{\longrightarrow}
\newcommand{\embbed}{\hookrightarrow}
\renewcommand{\bar}{\overline}

\newcommand{\id}{\operatorname{id}}
\newcommand{\im}{\operatorname{im}}

\newcommand{\normal}{\unlhd}

\renewcommand{\tilde}{\widetilde}
\renewcommand{\hat}{\widehat}
\newcommand{\del}{\partial}

\newcommand{\abs}[1]{\left| #1 \right|}
\newcommand{\norm}[1]{\left\lVert #1 \right\rVert}
\newcommand{\br}[1]{\left( #1 \right)}
\newcommand{\set}[1]{\left\{ #1 \right\}}
\newcommand{\cor}[1]{\left[ #1 \right]}
\newcommand{\gen}[1]{\left\langle #1 \right\rangle}
\newcommand{\matr}[1]{\begin{pmatrix}#1\end{pmatrix}}
\newcommand{\dv}[2]{\frac{\d#1}{\d#2}}

% Por si me equivoco al escribir
\newcommand{\rac}[2]{\frac{#1}{#2}}
\newcommand{\olon}{\colon}
\newcommand{\ubseteq}{\subseteq}
\newcommand{\ubset}{\subset}
\newcommand{\irc}{\circ}
\renewcommand{\o}{\to}

\newcommand{\ds}{\displaystyle}

\renewcommand*\contentsname{Contenidos}

\newcommand{\wea}[1]{\textcolor{BlueViolet}{\textbf{\emph{#1}}}}
\newcommand{\fecha}[1]{\begin{flushright} \underline{#1} \end{flushright}}

\newcommand{\dem}[1]{\textbf{Demostración (#1)}}


\title{MAT2504 -- Ecuaciones Diferenciales Ordinarias}
\author{Felipe Inostroza \\ Prof. Amal Taarabt}
\date{Segundo Semestre 2026}


\begin{document}

\maketitle
\tableofcontents

\newpage

\section{Introducción}
\subsection{Generalidades y Definiciones}
\fecha{Clase 1. Jueves 6 de Agosto}

Las leyes y los fenómenos del universo están escritos en el lenguaje matemático, y uno de los principales problemas de los científicos antiguos fué el estudio de los movimientos planetarios.

\begin{definition}
    Una \wea{ecuación diferencial ordinaria} (\wea{EDO}) de \wea{orden $k$} es una relación funcional entre la variable real $t \in \R$, una función $t \mapsto y(t)$ y sus derivadas $y', y'', \dots, y^{(k)}$ en $t$, definida por
    \[ F(t,y,y',y'',\dots,y^{(k)}) = 0, \tag{1} \]
    donde $y \colon J \to \R^m$ y $J \ubset \R$ es un intervalo abierto y 
    \[ F \colon J \times \R^m \times \dots \times \R^m \to \R^n \]
    es una función dada ($F \colon J \times \R^{(k+1)m} \to \R^n$).
    Generalmente $t$ se llama la \wea{variable independiente} e $y$ la \wea{variable dependiente}.
\end{definition}

\begin{example}\
    \begin{itemize}
        \item $f' + 2f = 10$ es una EDO, con la variable independiente dada implícitamente. La relación funcional es
        \[ F(x,f,f') \mapsto f' + 2f - 10. \]
        \item $f' + 2f = e^x \sin x$ es una EDO, con relación funcional
        \[ F(x,f,f') \mapsto f' + 2f - e^x \sin x. \]
    \end{itemize}
\end{example}

\begin{definition}
    Si $F$ es una función a valores reales, la EDO se dice \wea{ecuación escalar}.
\end{definition}

\begin{obs}
    En general $m = 1$, es decir $y \colon J \to \R$.
\end{obs}

\begin{example}
    \
    \begin{itemize}
        \item $ty' + 5y^2 = 5t^3$.
        \item $t^2 y'' - 2(y')^2 = 0$.
    \end{itemize}
\end{example}

\begin{notation}
$C^k(U,V)$ es el conjunto de las funciones $f \colon U \to V$ de clase $C^k$. Usamos también las notaciones
\[ y^{(n)}(t) = \dv{^n y}{t^n}. \]
\end{notation}

\begin{definition}
    Si $F$ no es independiente en $y^{(k)}$, entonces diremos que
    \[ y^{(k)} = f(t,y,y',\dots,y^{(k-1)}) \]
    es un sistema de $m \times m$ ecuaciones diferenciales ordinarias de orden $k$ con respecto a la función desconocida $y \colon J \to \R^m$ y  donde la función $f \colon J \times \R^{km} \to \R^m$. Escribimos
    \[ y_i^{(k)} = f_i(t,y,y',y^{(k-1)}), \quad \forall i \in \{1,\dots,m\} \]
    donde $y = (y_1,\dots,y_m)$ y $f = (f_1,\dots,f_m)$.
\end{definition}

\begin{example}
    \[ ty + (1+t^2) y' + \sin t y'' = 0. \]
    En este caso tenemos $F \colon (t,y,y',y'') \mapsto ty + (1+t^2) y' + \sin t y''$.
    Si consideramos momentáneamente valores de $t$ tales que $\sin t \neq 0$, entonces podemos reescribir esta EDO de la forma
    \[ y'' = - \frac{t}{\sin t} y - \frac{(1+t^2)}{\sin t} y'. \]
    Si tomamos $\ds f \colon (x,y,y') \mapsto - \frac{t}{\sin t} y - \frac{(1+t^2)}{\sin t} y'$, entonces pasamos el sistema de la forma $F(t,y,y',y'') = 0$ a $y'' = f(t,y,y')$ (forma estándar).
\end{example}

\begin{obs}
    \ \begin{itemize}
        \item En una EDO, la función desconocida $y$ depende de una sola variable independiente.
        \item Si la variable dependiente $y$ es una función de dos o más variables independientes, entonces aparecerán \wea{derivadas parciales} y hablaremos de EDP, ecuaciones a derivadas parciales. Por ejemplo, la temperatura $u(t,x)$ de una barra uniforme en el punto $x$ al tiempo $t$ satisface (junto con algunas condiciones adicionales) la siguiente EDP:
        \[ \frac{\del u}{\del t} = k \frac{\del^2 u}{\del x^2}, \]
        donde $k$ es la difusividad térmica de la barra (cst).
    \end{itemize}
\end{obs}

\begin{notation}
    Para derivadas parciales, escribimos
    \[ \frac{\del f}{\del x} = \del_x f. \]
\end{notation}

\begin{definition}
    Si $f \colon U \subset \R^n \to \R^m$, escribimos
    \[ \d f = \del_{x_1} f \d x_1 + \dots \del_{x_n} f \d x_n. \]
\end{definition}

\subsection{Solución de una ecuación diferencial}

Consideremos la EDO
\[ 2 f'' + f + (x^2 + 1) f = 0. \]
¿Cómo podemos resolverla? ¿Qué significa resolver una EDO? Queremos encontrar una función $f$ que satisfaga la relación funcional. También queremos saber en donde está definida, y claramente que sea de clase $C^2$ en su dominio.

\begin{definition}
    Una función $\phi \in C^k(I,\R^m)$ es \wea{solución} de la EDO (1) sobre $I \subseteq J$ si tenemos
    \[ \forall t \in I, \quad F(t,\phi(t),\phi'(t), \dots, \phi^{(k)}(t)) = 0. \]
    En este caso hablamos de una solución \wea{explícita}.
\end{definition}

\begin{example}
    La función $y(x) = 7 e^x + x^2 + 2x + 2$ es una solución de la EDO $y' = y - x^2$ en $I = \R$, ya que efectivamente
    \[ y'(x) = 7e^x + 2x + 2 = y(x) - 2. \]
\end{example}

\begin{definition}
    Una EDO de orden $k$ se dice de \wea{forma canónica} (\wea{normal}) si se escribe de la forma
    \[ y(k) = f(t,y,y',\dots,y^{(k-1)}), \tag{2} \]
    donde $f \colon I \times \R^{km} \to \R^m$ e $I \subset J$.
\end{definition}


\end{document}